\documentclass[11pt,a4paper]{article}

% === Essential Packages ===
\usepackage[utf8]{inputenc}
\usepackage[T1]{fontenc}
\usepackage{lmodern}

% === Math Packages ===
\usepackage{amsmath}
\usepackage{amssymb}
\usepackage{amsthm}
\usepackage{mathtools}

% === Tables and Figures ===
\usepackage{booktabs}
\usepackage{array}
\usepackage{graphicx}

% === Code Listings ===
\usepackage{listings}
\usepackage{xcolor}

% === Layout and Formatting ===
\usepackage[margin=1in]{geometry}
\usepackage{parskip}
\usepackage{enumitem}

% === Hyperlinks ===
\usepackage{hyperref}
\hypersetup{
    colorlinks=true,
    linkcolor=blue,
    citecolor=blue,
    urlcolor=blue
}

% === Algorithm Package ===
\usepackage[ruled,vlined,linesnumbered]{algorithm2e}

% === Code Listing Style ===
\definecolor{codegreen}{rgb}{0,0.6,0}
\definecolor{codegray}{rgb}{0.5,0.5,0.5}
\definecolor{codepurple}{rgb}{0.58,0,0.82}
\definecolor{backcolour}{rgb}{0.95,0.95,0.92}

\lstdefinestyle{mystyle}{
    backgroundcolor=\color{backcolour},
    commentstyle=\color{codegreen},
    keywordstyle=\color{magenta},
    numberstyle=\tiny\color{codegray},
    stringstyle=\color{codepurple},
    basicstyle=\ttfamily\footnotesize,
    breakatwhitespace=false,
    breaklines=true,
    captionpos=b,
    keepspaces=true,
    numbers=left,
    numbersep=5pt,
    showspaces=false,
    showstringspaces=false,
    showtabs=false,
    tabsize=2,
    frame=single
}
\lstset{style=mystyle}

% === Custom Commands ===
\newcommand{\ziptoken}[1]{\texttt{<ZIP\_#1>}}
\newcommand{\E}{\mathbb{E}}
\newcommand{\R}{\mathbb{R}}
\newcommand{\vocab}{\mathcal{V}}

% === Theorem Environments ===
\theoremstyle{definition}
\newtheorem{definition}{Definition}[section]
\newtheorem{property}{Property}[section]

% === Document Info ===
\title{\textbf{ZIP-RC: Zero-Overhead Introspection for\\Adaptive Test-Time Compute}\\[0.5em]
\large A Framework for LLM Self-Prediction of Success and Compute Cost}
\author{ZIP-RC Framework Documentation}
\date{\today}

\begin{document}

\maketitle

\begin{abstract}
ZIP-RC (Zero-Overhead Introspection for Adaptive Test-Time Compute) is a framework that enables Large Language Models (LLMs) to predict their own success probability and computational cost during inference---without adding any extra inference overhead. By extending the model vocabulary with special ``ZIP tokens,'' the model learns to jointly predict (1) whether the final answer will be correct (reward) and (2) how many more tokens remain until completion (length). This document provides a comprehensive technical explanation of the ZIP-RC architecture, mathematical formulations, code components, and practical applications.
\end{abstract}

\tableofcontents
\newpage

%==============================================================================
\section{Introduction}
%==============================================================================

Large Language Models have demonstrated remarkable capabilities in reasoning and problem-solving tasks. However, during generation, models typically operate without awareness of their own success probability or remaining computational cost. ZIP-RC addresses this limitation by enabling models to ``introspect'' at every generation step.

\subsection{Key Innovation}

The model learns to \textbf{jointly predict} two quantities at every generation step:
\begin{enumerate}
    \item \textbf{Reward}: Will the final answer be correct? (binary: wrong/correct)
    \item \textbf{Remaining Length}: How many more tokens until completion? (5 bins: 0--1, 2--3, 4--7, 8--15, 16+)
\end{enumerate}

This allows the model to make adaptive decisions during generation, such as when to stop, when to branch, and which paths to prune.

\subsection{Zero Overhead Property}

The critical advantage of ZIP-RC is that introspection comes at \textbf{zero additional computational cost}. The ZIP predictions are derived from the same forward pass used for text generation---no extra model calls are required.

%==============================================================================
\section{Architecture Overview}
%==============================================================================

ZIP-RC consists of two main components: a modified model architecture and a data generation pipeline.

\subsection{ZipRCModel: Core Model Wrapper}

The \texttt{ZipRCModel} class extends any causal language model with reserved ZIP tokens:
\begin{itemize}
    \item Adds $B_V \times B_T = 2 \times 5 = 10$ special tokens: \ziptoken{0} through \ziptoken{9}
    \item Each token represents one cell in the (reward\_bin, length\_bin) grid
    \item A single forward pass outputs both text logits AND introspection logits
\end{itemize}

\subsection{Data Pipeline}

The training data is generated through a four-stage process:
\begin{enumerate}
    \item \texttt{generate\_rollouts.py} $\rightarrow$ Generate completions from base model
    \item \texttt{score\_rollouts.py} $\rightarrow$ Score each completion (correct/incorrect)
    \item \texttt{build\_prefix\_dataset.py} $\rightarrow$ Create training examples from prefixes
    \item Training loop $\rightarrow$ Teach the model to predict (reward, length) at each prefix
\end{enumerate}

\subsection{Training Data Structure}

For each completion with $T$ tokens, we create $T$ training examples. Example $t$ predicts: (final\_reward, tokens\_remaining $= T - t$). This teaches the model to look ahead at every generation step.

\subsection{Pipeline Flow}

The overall pipeline can be visualized as:
\begin{center}
\begin{tabular}{ccccc}
\textbf{Prompts} & $\rightarrow$ & \textbf{Generate} & $\rightarrow$ & \textbf{Score} \\
(JSONL) & & Rollouts Using & & Each Answer \\
& & Base LM & & \\[0.5em]
& & $\downarrow$ & & \\[0.5em]
\textbf{Train} & $\leftarrow$ & \textbf{Build Dataset} & $\leftarrow$ & \\
ZipRC Model & & Prefix Examples & & \\
& & with Labels & &
\end{tabular}
\end{center}

%==============================================================================
\section{Mathematical Formulation}
%==============================================================================

This section presents the formal mathematical framework underlying ZIP-RC.

\subsection{Model Architecture}

\subsubsection{Vocabulary Extension}

The ZIP-RC model extends the original vocabulary with special tokens:
\begin{equation}
\vocab_{\text{ZIP}} = \vocab_{\text{original}} \cup \{\langle\text{ZIP}_0\rangle, \ldots, \langle\text{ZIP}_{B_V \cdot B_T - 1}\rangle\}
\end{equation}
where:
\begin{itemize}
    \item $B_V = 2$ (reward bins: wrong/correct)
    \item $B_T = 5$ (length bins: exponentially-spaced buckets)
    \item Total ZIP tokens: $B_V \times B_T = 10$
\end{itemize}

\subsubsection{Forward Pass}

Given an input sequence $\mathbf{x}_{1:t}$, the model produces logits:
\begin{equation}
\text{logits} = f_\theta(\mathbf{x}_{1:t}) \in \R^{|\vocab_{\text{ZIP}}|}
\end{equation}

These logits are split into two parts:
\begin{align}
\text{logits}_{\text{text}} &= \text{logits}[\,:\,|\vocab_{\text{original}}|] \in \R^{|\vocab_{\text{original}}|} \\
\text{logits}_{\text{ZIP}} &= \text{logits}[-B_V \cdot B_T\,:] \in \R^{B_V \times B_T}
\end{align}

\subsubsection{Generation Masking (Zero Overhead)}

During text generation, ZIP tokens are masked to ensure they are never sampled:
\begin{equation}
\text{logits}_{\text{gen}}[\text{ZIP indices}] = -\infty
\end{equation}

This ensures ZIP tokens are \textbf{never sampled} as text during generation, maintaining the model's language generation capabilities while enabling introspection.

\subsection{Joint Prediction}

\subsubsection{Grid Cell Mapping}

The ZIP tokens form a 2D grid where each cell corresponds to a (reward, length) pair:
\begin{equation}
\text{cell\_id}(v, \ell) = v \cdot B_T + \ell
\end{equation}
where:
\begin{itemize}
    \item $v \in \{0, 1\}$ is the reward bin (0 = wrong, 1 = correct)
    \item $\ell \in \{0, 1, 2, 3, 4\}$ is the length bin
\end{itemize}

\subsubsection{Length Binning Function}

The number of remaining tokens is discretized into bins:
\begin{equation}
\text{bin}(\text{tokens\_left}) = \begin{cases}
0 & \text{if tokens\_left} \in [0, 1] \\
1 & \text{if tokens\_left} \in [2, 3] \\
2 & \text{if tokens\_left} \in [4, 7] \\
3 & \text{if tokens\_left} \in [8, 15] \\
4 & \text{if tokens\_left} \geq 16
\end{cases}
\end{equation}

The bins are exponentially spaced to provide finer resolution for shorter remaining lengths.

\subsubsection{Joint Distribution}

The joint probability over reward and length is obtained via softmax:
\begin{equation}
p(v, \ell \mid \mathbf{x}_{1:t}) = \text{softmax}(\text{logits}_{\text{ZIP}})_{v \cdot B_T + \ell}
\end{equation}

\subsubsection{Marginal Distributions}

The marginal distributions can be computed by summing over the joint:
\begin{align}
p(v \mid \mathbf{x}_{1:t}) &= \sum_{\ell=0}^{B_T-1} p(v, \ell \mid \mathbf{x}_{1:t}) \quad \text{(reward probability)} \\
p(\ell \mid \mathbf{x}_{1:t}) &= \sum_{v=0}^{B_V-1} p(v, \ell \mid \mathbf{x}_{1:t}) \quad \text{(length probability)}
\end{align}

\subsubsection{Expected Values}

Expected reward and length can be computed as:
\begin{align}
\E[v] &= \sum_{v=0}^{B_V-1} v \cdot p(v \mid \mathbf{x}_{1:t}) \\
\E[\ell] &= \sum_{\ell=0}^{B_T-1} \ell \cdot p(\ell \mid \mathbf{x}_{1:t})
\end{align}

\subsection{Training Objective}

For a completion with $T$ tokens, we create $T$ training examples.

\subsubsection{Label Construction}

At prefix position $t$:
\begin{align}
\text{tokens\_left} &= T - t \\
v^* &= \mathbb{1}[\text{completion is correct}] \\
\ell^* &= \text{bin}(\text{tokens\_left}) \\
y^* &= v^* \cdot B_T + \ell^* \quad \text{(joint label)}
\end{align}

\subsubsection{Loss Function}

The training uses standard cross-entropy loss over the 10 joint classes:
\begin{equation}
\mathcal{L} = -\log p(v^*, \ell^* \mid \mathbf{x}_{1:t}) = -\log \text{softmax}(\text{logits}_{\text{ZIP}})_{y^*}
\end{equation}

\subsection{Key Properties}

\begin{property}[Zero Overhead]
The computational cost of ZIP-RC inference equals that of the base model:
\begin{equation}
\text{Time}(\text{ZIP-RC forward pass}) = \text{Time}(\text{base LM forward pass})
\end{equation}
The ZIP logits are computed from the same hidden states---no extra forward passes needed.
\end{property}

\begin{property}[Introspection at Every Step]
At any generation step $t$, the model can query:
\begin{align}
p(\text{success} \mid \text{prefix}_{1:t}) &= p(v=1 \mid \mathbf{x}_{1:t}) \\
\E[\text{remaining tokens} \mid \text{prefix}_{1:t}] &\approx \text{unbin}(\E[\ell])
\end{align}
\end{property}

%==============================================================================
\section{Code Components Explained}
%==============================================================================

This section provides detailed explanations of each code component in the ZIP-RC implementation.

\subsection{ZipRCModel}

The core model class in \texttt{zip\_rc\_model.py}.

\subsubsection{Initialization}

The model initialization performs the following steps:
\begin{enumerate}
    \item Loads base causal LM (default: Meta-Llama-3.1-8B-Instruct)
    \item Adds 10 ZIP tokens to vocabulary: \ziptoken{0} through \ziptoken{9}
    \item Resizes embedding matrix to accommodate new tokens
    \item Optionally freezes backbone, training only the LM head
\end{enumerate}

\subsubsection{Forward Pass}

The forward pass consists of three stages:
\begin{enumerate}
    \item Run base model forward: \texttt{logits = model(input\_ids)}
    \item Split logits into text vs ZIP:
    \begin{itemize}
        \item \texttt{normal\_logits = logits[:, :original\_vocab\_size]}
        \item \texttt{zip\_logits = logits[:, -10:]}
    \end{itemize}
    \item Create generation logits with ZIP tokens masked to $-\infty$
    \item Return all three: \texttt{token\_logits}, \texttt{generation\_logits}, \texttt{zip\_logits}
\end{enumerate}

\subsubsection{Prediction}

The prediction method:
\begin{enumerate}
    \item Applies softmax to ZIP logits $\rightarrow$ joint distribution $p(v, \ell)$
    \item Reshapes into $2 \times 5$ grid
    \item Sums over dimensions to get marginals
    \item Computes expected reward and expected length
\end{enumerate}

\subsection{generate\_rollouts.py}

\textbf{Purpose:} Generate completions from base model (no ZIP tokens yet).

\textbf{Process:}
\begin{enumerate}
    \item Load prompts from \texttt{data/prompts.jsonl}
    \item Format as ``Q: \{prompt\}\textbackslash nA:''
    \item Generate greedy completions (max 256 tokens)
    \item Save to \texttt{data/rollouts.jsonl} with: prompt, answer (ground truth), completion, completion\_ids
\end{enumerate}

\subsection{score\_rollouts.py}

\textbf{Purpose:} Evaluate each completion for correctness.

\textbf{Scoring Logic:}
\begin{enumerate}
    \item Extract last number from completion text
    \item Compare to ground truth answer
    \item Assign reward: 1 if match, 0 otherwise
    \item Add \texttt{predicted\_answer} and \texttt{reward} fields
\end{enumerate}

\textbf{Output:} \texttt{data/scored\_rollouts.jsonl}

\subsection{build\_prefix\_dataset.py}

\textbf{Purpose:} Create training examples for ZIP-RC.

\textbf{Key Algorithm:}

\begin{algorithm}[H]
\SetAlgoLined
\KwIn{Scored rollouts with completions and rewards}
\KwOut{Prefix dataset with joint labels}
\ForEach{scored rollout}{
    $T \leftarrow$ length of completion\;
    \For{$t \leftarrow 1$ \KwTo $T$}{
        prefix $\leftarrow$ prompt $+$ completion$[{:}t]$\;
        tokens\_left $\leftarrow T - t$\;
        length\_bin $\leftarrow$ bin(tokens\_left)\;
        joint\_label $\leftarrow$ reward\_bin $\times 5 +$ length\_bin\;
        Save: \{input\_ids, attention\_mask, joint\_label\}\;
    }
}
\caption{Build Prefix Dataset}
\end{algorithm}

\textbf{Example:}
\begin{itemize}
    \item Completion has $T=20$ tokens, reward$=1$ (correct)
    \item At $t=15$: tokens\_left$=5 \rightarrow$ length\_bin$=2$ (bucket 4--7)
    \item joint\_label $= 1 \times 5 + 2 = 7$ (ZIP token 7)
\end{itemize}

\textbf{Output:} \texttt{data/prefix\_dataset.jsonl} with thousands of prefix examples.

\subsection{Training Loop}

The training follows standard supervised learning:

\begin{lstlisting}[language=Python, caption=Training Loop Pseudocode]
model = ZipRCModel()
for batch in dataset:
    outputs = model(input_ids, attention_mask)
    zip_logits = outputs.zip_logits[:, -1, :]  # last position
    loss = CrossEntropyLoss(zip_logits, joint_labels)
    loss.backward()
\end{lstlisting}

The model learns to predict the correct ZIP token at each prefix position.

%==============================================================================
\section{Complete Pipeline Example}
%==============================================================================

This section walks through a complete example using the prompt ``What is 15 + 27?''

\subsection{Stage 1: Generate Rollout}

\begin{verbatim}
Prompt: "What is 15 + 27?"
Ground truth: "42"
Base model generates: "15 + 27 = 42"
Completion tokens: [220, 16, 20, 488, 220, 17, 22, 284, 220, 19, 17]
                   (11 tokens)
\end{verbatim}

\subsection{Stage 2: Score}

\begin{verbatim}
Extract last number: "42"
Compare to ground truth: "42"
Reward: 1 (correct)
\end{verbatim}

\subsection{Stage 3: Build Prefix Dataset}

We create 11 training examples, one for each prefix:

\begin{table}[h]
\centering
\caption{Training Examples Generated from ``15 + 27 = 42''}
\begin{tabular}{@{}cllccc@{}}
\toprule
$t$ & Prefix & Tokens Left & Length Bin & Joint Label & ZIP Token \\
\midrule
1 & ``15'' & 10 & 3 & $1 \times 5 + 3 = 8$ & \ziptoken{8} \\
2 & ``15 +'' & 9 & 3 & 8 & \ziptoken{8} \\
3 & ``15 + '' & 8 & 3 & 8 & \ziptoken{8} \\
4 & ``15 + 27'' & 7 & 2 & $1 \times 5 + 2 = 7$ & \ziptoken{7} \\
5 & ``15 + 27 '' & 6 & 2 & 7 & \ziptoken{7} \\
6 & ``15 + 27 ='' & 5 & 2 & 7 & \ziptoken{7} \\
7 & ``15 + 27 = '' & 4 & 2 & 7 & \ziptoken{7} \\
8 & ``15 + 27 = 4'' & 3 & 1 & $1 \times 5 + 1 = 6$ & \ziptoken{6} \\
9 & ``15 + 27 = 42'' & 2 & 1 & 6 & \ziptoken{6} \\
10 & ``15 + 27 = 42'' & 1 & 0 & $1 \times 5 + 0 = 5$ & \ziptoken{5} \\
11 & ``15 + 27 = 42'' (end) & 0 & 0 & 5 & \ziptoken{5} \\
\bottomrule
\end{tabular}
\end{table}

\subsection{Stage 4: Training}

For each training example:
\begin{enumerate}
    \item Feed prefix into model
    \item Model outputs 10 ZIP logits
    \item Loss $=$ CrossEntropy(zip\_logits, target\_label)
    \item Model learns: ``After `15 + 27 =', I should output \ziptoken{7} (correct, 4--7 tokens left)''
\end{enumerate}

\subsection{Stage 5: Inference (After Training)}

\begin{lstlisting}[language=Python, caption=Inference Example]
model = ZipRCModel.load_trained()
prompt = "What is 15 + 27?"
prefix = "15 + 27 = 4"

reward, length = model.predict_zip(prefix)
# reward ~ 0.95 (high confidence of correctness)
# length ~ 0.8 (likely 0-3 tokens remaining)
\end{lstlisting}

The model now knows mid-generation whether it's on track!

%==============================================================================
\section{Use Cases and Applications}
%==============================================================================

ZIP-RC enables several practical applications for adaptive inference.

\subsection{Adaptive Early Stopping}

\begin{lstlisting}[language=Python, caption=Early Stopping Strategy]
if p(success | prefix) < 0.3:
    stop_generation()
    try_different_approach()
\end{lstlisting}

When the model detects low success probability, it can abort the current generation path and try an alternative approach.

\subsection{Beam Search Pruning}

\begin{lstlisting}[language=Python, caption=Beam Pruning Strategy]
for beam in beams:
    score = p(success) - cost_weight * expected_length
    beam.score = score
beams = sorted(beams, key=lambda b: b.score, reverse=True)[:k]
\end{lstlisting}

Beams can be scored and pruned based on both success probability and expected remaining length.

\subsection{Best-of-N Selection}

\begin{lstlisting}[language=Python, caption=Best-of-N Selection]
completions = generate_n_completions(prompt, n=N)
best_idx = argmax([p(success | c) for c in completions])
return completions[best_idx]
\end{lstlisting}

When generating multiple candidate completions, select the one with highest predicted success probability.

\subsection{Dynamic Compute Allocation}

\begin{lstlisting}[language=Python, caption=Dynamic Compute Allocation]
if p(success) > 0.8 and expected_length < 5:
    # Efficient path: continue
    continue_generation()
elif p(success) < 0.2 and expected_length > 20:
    # Waste of compute: abort early
    abort_and_retry()
\end{lstlisting}

Allocate compute resources based on the expected utility of continuing generation.

\subsection{Uncertainty Estimation}

\begin{lstlisting}[language=Python, caption=Uncertainty-Based Branching]
entropy = -sum(p(v,l) * log(p(v,l)) for v,l in grid)
if entropy > threshold:
    # Model is uncertain: consider multiple paths
    branch_generation()
\end{lstlisting}

High entropy in the joint distribution indicates model uncertainty, suggesting multiple generation paths should be explored.

\subsection{Summary of Key Benefits}

\begin{description}
    \item[Zero Overhead:] No extra forward passes---introspection is ``free''
    \item[Real-time:] Predictions available at every token
    \item[Joint Modeling:] Captures correlation between success and length
    \item[Generalizable:] Works with any causal LM backbone
\end{description}

%==============================================================================
\section{File Structure Summary}
%==============================================================================

The ZIP-RC codebase is organized as follows:

\begin{verbatim}
zip-rc/
├── zip_rc_model.py          # Core model: adds ZIP tokens, splits logits
├── generate_rollouts.py     # Step 1: Generate completions from base LM
├── score_rollouts.py        # Step 2: Evaluate correctness
├── build_prefix_dataset.py  # Step 3: Create prefix training examples
├── build_rollout_labels.py  # Alternative labeling script (more general)
├── data/
│   ├── prompts.jsonl        # Input: math questions + answers
│   ├── rollouts.jsonl       # Output: generated completions
│   ├── scored_rollouts.jsonl # Output: completions + rewards
│   └── prefix_dataset.jsonl # Output: training examples
└── README.md                # Citation and overview
\end{verbatim}

\subsection{Execution Order}

\begin{enumerate}
    \item \texttt{python generate\_rollouts.py} $\rightarrow$ creates \texttt{rollouts.jsonl}
    \item \texttt{python score\_rollouts.py} $\rightarrow$ creates \texttt{scored\_rollouts.jsonl}
    \item \texttt{python build\_prefix\_dataset.py} $\rightarrow$ creates \texttt{prefix\_dataset.jsonl}
    \item Train ZipRCModel on \texttt{prefix\_dataset.jsonl} (training script not included)
    \item Use trained model for adaptive inference
\end{enumerate}

\subsection{Model Files}

\begin{description}
    \item[Base model:] \texttt{meta-llama/Meta-Llama-3.1-8B-Instruct} (from HuggingFace)
    \item[Trained ZIP-RC checkpoint:] (train yourself using prefix dataset)
\end{description}

%==============================================================================
\section{Conclusion}
%==============================================================================

ZIP-RC provides an elegant solution to the problem of LLM self-awareness during generation. By extending the vocabulary with special tokens and training on prefix-labeled data, the model gains the ability to predict its own success probability and remaining generation length at zero additional inference cost.

The framework's key contributions are:
\begin{enumerate}
    \item A simple architectural modification that adds introspection capabilities to any causal LM
    \item A systematic data generation pipeline for creating training examples
    \item Mathematical foundations for joint reward-length prediction
    \item Practical applications for adaptive test-time compute allocation
\end{enumerate}

Future work may explore:
\begin{itemize}
    \item More fine-grained length binning for longer generations
    \item Multi-dimensional reward signals beyond binary correctness
    \item Integration with tree-search methods like MCTS
    \item Application to domains beyond mathematical reasoning
\end{itemize}

\end{document}
